\documentclass[12pt]{article}
\usepackage[a4paper, top=2.5cm, bottom=2.5cm, left=1.5cm, right=1.5cm]{geometry}
\usepackage{amsmath, amsfonts, amssymb, mathtools}
\usepackage{fancyhdr, setspace, parskip}
\usepackage{graphicx, caption, subfig, array, multirow}
\usepackage{hyperref, enumitem, cancel}
\usepackage[T1]{fontenc}
\usepackage{tgtermes}
\usepackage[dvipsnames]{xcolor}
\usepackage{tocloft}
\usepackage{titlesec}
\usepackage{lipsum}  

\definecolor{DarkBlue}{RGB}{10, 0, 80}

% Hyperlink setup
\hypersetup{
    colorlinks=true,
    linkcolor=DarkBlue,
    filecolor=BrickRed,      
    urlcolor=RoyalBlue,
}


% Header and footer customization
\fancyhead{}
\fancyhead[L]{
{\fontfamily{lmss}{\color{DarkBlue}
\textbf{\leftmark}
}}
}
\fancyhead[R]{
{\fontfamily{ppl}\selectfont {\color{DarkBlue}
{Deep RL [Spring 2025]}
}}
}

\fancyfoot{}
\fancyfoot[C]{
{\fontfamily{lmss}{\color{BrickRed}
\textbf{\thepage}
}}
}

\renewcommand{\sectionmark}[1]{ \markboth{\thesection\quad #1}{} }

\renewcommand{\headrule}{{\color{BrickRed}\hrule width\headwidth height 0.5pt}}
\renewcommand{\footrulewidth}{0pt}


% Table of Contents customizations
\renewcommand{\cftsecafterpnum}{\vskip6pt}
\renewcommand{\cftsubsecafterpnum}{\vskip3pt}
\renewcommand{\cftsubsubsecafterpnum}{\vskip3pt}
\renewcommand{\cftsecfont}{\sffamily\large}
\renewcommand{\cftsubsecfont}{\sffamily}
\renewcommand{\cftsubsubsecfont}{\sffamily}
% \renewcommand{\cftsecdotsep}{1}
\renewcommand{\cftsubsecdotsep}{1}
\renewcommand{\cftsubsubsecdotsep}{1}


% Section title styles
\titleformat*{\section}{\LARGE\bfseries\color{DarkBlue}}
\titleformat*{\subsection}{\Large\bfseries\color{DarkBlue}}
\titleformat*{\subsubsection}{\large\bfseries\color{DarkBlue}}

\definecolor{light-gray}{gray}{0.95}
\newcommand{\code}[1]{\colorbox{light-gray}{\texttt{#1}}}

% Start of the document
\pagestyle{fancy}

%%%%%%%%%%%%%%%%%%%%%%%%%%%%%%%%%%%%%%%%%%%%%%%%%

\begin{document}

\pagenumbering{gobble}
\thispagestyle{plain}

\begin{center}

\vspace*{-1.5cm}
\begin{figure}[!h]
    \centering
    \includegraphics[width=0.7\linewidth]{figs/cover-std.png}
\end{figure}

{
\fontfamily{ppl}

{\color{DarkBlue} {\fontsize{30}{50} \textbf{
Deep Reinforcement Learning
}}}

{\color{DarkBlue} {\Large
Professor Mohammad Hossein Rohban
}}
}


\vspace{20pt}

{
\fontfamily{lmss}


{\color{RedOrange}
{\Large
Homework 1:
}\\
}
{\color{BrickRed}
\rule{12cm}{0.5pt}

{\Huge
Introduction to RL
}
\rule{12cm}{0.5pt}
}

\vspace{10pt}

{\color{RoyalPurple} { \small By:} } \\
\vspace{10pt}

{\color{Blue} { \LARGE [Full Name] } } \\
\vspace{5pt}
{\color{RoyalBlue} { \Large [Student Number] } }


\vspace*{\fill}
\begin{center}
\begin{tabular}{ccc}
    \includegraphics[width=0.14\linewidth]{figs/sharif-logo.png} & \includegraphics[width=0.14\linewidth]{figs/riml-logo.png} & \includegraphics[width=0.14\linewidth]{figs/dlr-logo.png} \\
\end{tabular}
\end{center}

\vspace*{-.25cm}

{\color{YellowOrange} {
\rule{10cm}{0.5pt} \\
\vspace{2pt}
\large Spring 2025}
}}
\vspace*{-1cm}

\end{center}

%%%%%%%%%%%%%%%%%%%%%%%%%%%%%%%%%%%%%%%%%%%%%%%%%

\newpage
\pagenumbering{gobble}
\thispagestyle{plain}
{\fontfamily{lmss}\selectfont {\color{BrickRed} \textbf{\tableofcontents} }}

{\fontfamily{lmss}\selectfont {\color{DarkBlue}

\subsection*{Grading}

The grading will be based on the following criteria, with a total of 100 points:

\[
\begin{array}{|l|l|}
\hline
\textbf{Task} & \textbf{Points} \\
\hline
\text{Task 1: Solving Predefined Environments} & 45 \\
\text{Task 2: Creating Custom Environments} & 45 \\

\hline
\text{Clarity and Quality of Code} & 5 \\
\text{Clarity and Quality of Report} & 5 \\
\hline
\text{Bonus 1: Writing a wrapper for a known env} & 10 \\
\text{Bonus 2: Implementing pygame env } & 20 \\
\text{Bonus 3: Writing your report in Latex } & 10 \\
\hline
\end{array}
\]

\textbf{Notes:}
\begin{itemize}
    \item Include well-commented code and relevant plots in your notebook.
    \item Clearly present all comparisons and analyses in your report.
    \item Ensure reproducibility by specifying all dependencies and configurations.
\end{itemize}

}

%%%%%%%%%%%%%%%%%%%%%%%%%%%%%%%%%%%%%%%%%%%%%%%%%

\newpage
\pagenumbering{arabic}

{\fontfamily{lmss}\selectfont {\color{DarkBlue}

\section{Task 1: Solving Predefined Environments [45-points]}

In this task, we explore reinforcement learning by training agents on predefined environments using Stable-Baselines3. We implement and compare different RL algorithms including A2C, PPO, DQN, and Q-Learning across multiple environments.

\subsection{Environment Selection and Implementation}

We selected four environments for comprehensive analysis:

\begin{itemize}
    \item \textbf{CartPole-v1}: A classic control problem where the agent must balance a pole on a cart
    \item \textbf{FrozenLake-v1}: A gridworld navigation problem with slippery surfaces
    \item \textbf{Taxi-v3}: A more complex navigation problem requiring pickup and dropoff
    \item \textbf{FlappyBird-v0}: A custom environment simulating the popular Flappy Bird game
\end{itemize}

\subsection{CartPole Environment Analysis}

\subsubsection{Environment Setup and Training}

The CartPole environment was trained using A2C and PPO algorithms:

\begin{verbatim}
# A2C Training
from stable_baselines3 import A2C, PPO, DQN
import gymnasium as gym

env = gym.make("CartPole-v1")
model = A2C("MlpPolicy", env, verbose=1, 
           tensorboard_log="./a2c_CartPole_tensorboard/")
model.learn(total_timesteps=500_000)
model.save("a2c_cartpole")

# PPO Training
model = PPO("MlpPolicy", env, verbose=1, 
           tensorboard_log="./PPO_CartPole_tensorboard/")
model.learn(total_timesteps=500_000)
model.save("ppo_cartpole")
\end{verbatim}

\subsubsection{Hyperparameter Analysis}

We conducted hyperparameter tuning experiments with different learning rates and discount factors:

\begin{verbatim}
learning_rates = [0.0001, 0.01]  
gammas = [0.95, 0.99]

for lr in learning_rates:
    for gamma in gammas:
        model = PPO("MlpPolicy", env, verbose=0,
                   learning_rate=lr, gamma=gamma,
                   tensorboard_log=f'./PPO_CartPole_tensorboard/ppo_lr{lr}_gamma{gamma}')
        model.learn(total_timesteps=200_000)
\end{verbatim}

\subsubsection{Reward Wrapping}

We implemented a reward wrapper to modify the reward function:

\begin{verbatim}
from gymnasium import RewardWrapper

class DoubledReward(RewardWrapper):
    def reward(self, reward):
        return reward * 2

wrapped_env = DoubledReward(env)
model = PPO("MlpPolicy", wrapped_env, verbose=0, 
           tensorboard_log="./PPO_CartPole_tensorboard/reward_wrapped_env_ppo")
model.learn(total_timesteps=500_000)
\end{verbatim}

\subsection{FrozenLake Environment Analysis}

\subsubsection{Algorithm Comparison}

For FrozenLake, we compared DQN and A2C algorithms:

\begin{verbatim}
env = gym.make("FrozenLake-v1", is_slippery=True)

# DQN Training
model = DQN("MlpPolicy", env, verbose=0, 
           tensorboard_log="./DQN_FrozenLake_tensorboard/")
model.learn(total_timesteps=500_000)

# A2C Training  
model = A2C("MlpPolicy", env, verbose=0, 
           tensorboard_log="./A2C_FrozenLake_tensorboard/")
model.learn(total_timesteps=500_000)
\end{verbatim}

\subsubsection{Custom Reward Function}

We implemented a modified reward function for FrozenLake:

\begin{verbatim}
class ModifiedFrozenLakeReward(RewardWrapper):
    def reward(self, reward):
        if reward == 0:            
            return -0.01
        elif reward == 1:           
            return 2
        else:
            return reward

wrapped_env = ModifiedFrozenLakeReward(env)
model = DQN("MlpPolicy", wrapped_env, verbose=0, 
           tensorboard_log="./DQN_FrozenLake_tensorboard/reward_wrapped_env_DQN")
model.learn(total_timesteps=500_000)
\end{verbatim}

\subsection{Taxi Environment Analysis}

\subsubsection{Q-Learning Implementation}

For the Taxi environment, we implemented a custom Q-Learning algorithm as it proved more effective than DQN and PPO:

\begin{verbatim}
def qlearning_train(taxi_env, learning_rate, discount_factor):
    writer = SummaryWriter(f"Qlearning_Taxi/taxi_experiment_lr_{learning_rate}_{discount_factor}")
    
    num_states = taxi_env.observation_space.n
    num_actions = taxi_env.action_space.n
    q_values = np.zeros((num_states, num_actions))
    
    alpha = learning_rate
    gamma = discount_factor
    eps = 1.0
    decay = 0.005
    episodes = 5000
    max_steps_per_ep = 99
    
    episode_rewards = []
    
    for ep in range(episodes):
        state, _ = taxi_env.reset()
        total_reward = 0
    
        for step in range(max_steps_per_ep):
            if random.random() < eps:
                chosen_action = taxi_env.action_space.sample()
            else:
                chosen_action = np.argmax(q_values[state, :])
            
            next_state, reward, done, _, _ = taxi_env.step(chosen_action)
            
            q_values[state, chosen_action] += alpha * (
                reward + gamma * np.max(q_values[next_state, :]) - q_values[state, chosen_action]
            )
    
            state = next_state
            total_reward += reward
    
            if done:
                break
    
        eps = 1 / (1 + decay * ep)
        episode_rewards.append(total_reward)
        
        avg_reward = np.mean(episode_rewards[-100:]) if ep >= 100 else np.mean(episode_rewards)
        writer.add_scalar("Mean_Episode_Reward", avg_reward, ep)
    
    return q_values
\end{verbatim}

\subsubsection{Agent Visualization}

We implemented a simulation function to visualize the trained agent:

\begin{verbatim}
def simulate_agent(environment, q_matrix, num_runs=5, max_steps_per_run=100):
    for run in range(num_runs):
        current_state, _ = environment.reset()
        finished = False
        
        for step in range(max_steps_per_run):            
            print(environment.render())
            sleep(0.5)  
            
            chosen_action = np.argmax(q_matrix[current_state, :])
            next_state, reward, finished, _, _ = environment.step(chosen_action)
            current_state = next_state

            if finished:
                print(f"Run completed in {step + 1} steps.\n")
                break
\end{verbatim}

\subsection{FlappyBird Environment Analysis}

\subsubsection{Environment Setup}

We used a third-party FlappyBird implementation:

\begin{verbatim}
!pip install flappy-bird-gymnasium

import flappy_bird_gymnasium
env = gymnasium.make("FlappyBird-v0", render_mode="rgb_array", use_lidar=True)

# Environment properties
print(f"Action Space: {env.action_space}")  # Discrete(2)
print(f"Observation Space: {env.observation_space}")  # Box(0.0, 1.0, (180,), float64)
\end{verbatim}

\subsubsection{Algorithm Training}

We trained A2C and PPO agents on FlappyBird:

\begin{verbatim}
# A2C Training
model = A2C("MlpPolicy", env, verbose=0, 
           tensorboard_log="./A2C_FlappyBird_tensorboard/")
model.learn(total_timesteps=5_000_000)

# PPO Training
model = PPO("MlpPolicy", env, verbose=0, 
           tensorboard_log="./PPO_mlp_FlappyBird_tensorboard/")
model.learn(total_timesteps=5_000_000)
\end{verbatim}

\subsection{Results and Analysis}

\subsubsection{Performance Comparison}

Through TensorBoard logging, we observed the following performance characteristics:

\begin{itemize}
    \item \textbf{CartPole}: Both A2C and PPO achieved stable performance, with PPO showing more consistent learning curves
    \item \textbf{FrozenLake}: DQN outperformed A2C due to the discrete nature of the environment
    \item \textbf{Taxi}: Q-Learning proved most effective, achieving optimal solutions in fewer episodes
    \item \textbf{FlappyBird}: Both A2C and PPO struggled with the high-dimensional observation space, requiring extensive training
\end{itemize}

\subsubsection{Hyperparameter Impact}

Our hyperparameter analysis revealed:

\begin{itemize}
    \item Higher learning rates (0.01) led to faster initial learning but less stable convergence
    \item Lower learning rates (0.0001) provided more stable learning but slower convergence
    \item Higher discount factors (0.99) were beneficial for long-term planning tasks
    \item Lower discount factors (0.95) worked better for immediate reward tasks
\end{itemize}

\subsubsection{Reward Wrapping Effects}

Custom reward functions significantly impacted learning:

\begin{itemize}
    \item Doubled rewards in CartPole accelerated learning
    \item Modified FrozenLake rewards improved exploration by penalizing non-progress
    \item Custom Taxi rewards with higher penalties for illegal actions improved policy quality
\end{itemize}

}}

%%%%%%%%%%%%%%%%%%%%%%%%%%%%%%%%%%%%%%%%%%%%%%%%%

\newpage

{\fontfamily{lmss}\selectfont {\color{DarkBlue}

\section{Task 2: Creating Custom Environments [45-points]}

In this task, we design and implement a custom 4×4 GridWorld environment following the Markov Decision Process (MDP) framework. We model the environment as a complete MDP, implement it using the Gymnasium API, and test it with various reinforcement learning algorithms.

\subsection{MDP Formulation}

\subsubsection{State Space (S)}

The state space consists of all 4×4 = 16 cells in the grid. Each state is represented as a tuple (row, column):

\begin{itemize}
    \item Start state: (0,0) - top-left corner marked as 'S'
    \item Goal state: (3,3) - bottom-right corner marked as 'G'  
    \item Hole states: (1,1) and (2,2) - marked as 'H'
    \item Regular states: remaining 10 cells
\end{itemize}

The state space is: $S = \{(0,0), (0,1), ..., (3,2), (3,3)\}$ with 16 possible states.

\subsubsection{Action Space (A)}

The agent can move in four cardinal directions:

\begin{itemize}
    \item Action 0: Move up
    \item Action 1: Move down  
    \item Action 2: Move left
    \item Action 3: Move right
\end{itemize}

The action space is: $A = \{0, 1, 2, 3\}$ with 4 possible actions.

\subsubsection{Reward Function (R)}

The reward function is defined as:

\begin{itemize}
    \item Reaching the goal (3,3): $+10$ reward
    \item Falling into a hole (1,1) or (2,2): $-1$ reward
    \item Moving to any other cell: $0$ reward
\end{itemize}

\subsubsection{Transition Probability (P)}

We assume a deterministic environment where actions always result in the intended movement, except when hitting boundaries (where the agent stays in place).

\subsection{Environment Implementation}

\subsubsection{Gymnasium API Implementation}

We implemented the GridWorld environment following the Gymnasium API:

\begin{verbatim}
import gymnasium as gym
from gymnasium import spaces
import numpy as np

class GridWorldEnv(gym.Env):
    def __init__(self):       
        self.grid_size = 4
        self.start = (0, 0)  
        self.goal = (3, 3)  
        self.holes = [(1, 1), (2, 2)]  
        self.state = self.start
        
        self.action_space = spaces.Discrete(4)  # 0: up, 1: down, 2: left, 3: right
        self.observation_space = spaces.Discrete(16)  # from 0 to 15

    def get_observation(self):
        """Convert internal state (i, j) to integer observation using state = 4*i + j."""
        i, j = self.state
        return 4 * i + j

    def step(self, action):        
        row, col = self.state
        
        # Calculate next state based on action
        if action == 0:  # up
            next_row = max(row - 1, 0)
            next_col = col
        elif action == 1:  # down
            next_row = min(row + 1, 3)
            next_col = col
        elif action == 2:  # left
            next_row = row
            next_col = max(col - 1, 0)
        elif action == 3:  # right
            next_row = row
            next_col = min(col + 1, 3)

        next_state = (next_row, next_col)
        
        # Determine rewards and termination
        if next_state in self.holes:
            reward = -1
            done = True
        elif next_state == self.goal:
            reward = 10
            done = True
        else:
            reward = 0
            done = False
        
        self.state = next_state
        return self.get_observation(), reward, done, False, {}

    def render(self):
        """Render the grid world with the agent's position."""
        grid = [['.' for _ in range(4)] for _ in range(4)]
        
        grid[0][0] = 'S'               # Start
        grid[3][3] = 'G'               # Goal
        for hole in self.holes:        # Holes
            grid[hole[0]][hole[1]] = 'H'
        
        agent_row, agent_col = self.state
        if grid[agent_row][agent_col] == 'S':
            grid[agent_row][agent_col] = 'A'  # Agent on start
        elif grid[agent_row][agent_col] == 'G':
            grid[agent_row][agent_col] = 'A'  # Agent on goal
        elif grid[agent_row][agent_col] == 'H':
            grid[agent_row][agent_col] = 'A(H)'  # Agent on hole
        else:
            grid[agent_row][agent_col] = 'A'     # Agent on empty cell
        
        for row in grid:
            print('    '.join(row))
        print()

    def reset(self):
        """Reset the environment to the start state."""
        self.state = self.start
        return self.get_observation(), {}
\end{verbatim}

\subsection{Q-Learning Implementation}

\subsubsection{Q-Learning Algorithm}

We implemented a custom Q-Learning algorithm specifically for the GridWorld environment:

\begin{verbatim}
def qlearning_train_gridworld(env, learning_rate, discount_factor):
    writer = SummaryWriter(f"Qlearning_GridWorld/experiment_lr_{learning_rate}_{discount_factor}")
    
    num_states = env.observation_space.n
    num_actions = env.action_space.n
    q_values = np.zeros((num_states, num_actions))
    
    alpha = learning_rate
    gamma = discount_factor
    eps = 1.0
    decay = 0.005
    episodes = 5000
    max_steps_per_ep = 99
    
    episode_rewards = []
    
    for ep in range(episodes):
        state, _ = env.reset()
        total_reward = 0
    
        for step in range(max_steps_per_ep):
            # Epsilon-greedy action selection
            if random.random() < eps:
                chosen_action = env.action_space.sample()
            else:
                chosen_action = np.argmax(q_values[state, :])
            
            next_state, reward, done, _, _ = env.step(chosen_action)
            
            # Q-value update
            q_values[state, chosen_action] += alpha * (
                reward + gamma * np.max(q_values[next_state, :]) - q_values[state, chosen_action]
            )
    
            state = next_state
            total_reward += reward
    
            if done:
                break
    
        # Epsilon decay
        eps = 1 / (1 + decay * ep)
        episode_rewards.append(total_reward)
        
        avg_reward = np.mean(episode_rewards[-100:]) if ep >= 100 else np.mean(episode_rewards)
        writer.add_scalar("Mean_Episode_Reward", avg_reward, ep)
    
    return q_values
\end{verbatim}

\subsubsection{Agent Simulation}

We implemented a visualization function to observe the trained agent's behavior:

\begin{verbatim}
def simulate_gridworld_agent(environment, q_matrix, max_steps_per_run=100):
    current_state, _ = environment.reset()
    finished = False    
    
    for step in range(max_steps_per_run):            
        environment.render()
        sleep(0.5)  
        
        chosen_action = np.argmax(q_matrix[current_state, :])
        next_state, reward, finished, _, _ = environment.step(chosen_action)
        current_state = next_state

        if finished:
            environment.render()
            print(f"Run completed in {step + 1} steps.\n")
            break
\end{verbatim}

\subsection{Algorithm Comparison}

\subsubsection{Q-Learning Performance}

We trained Q-Learning agents with different hyperparameters:

\begin{verbatim}
env = GridWorldEnv()
q_values = qlearning_train_gridworld(env=env, learning_rate=0.9, discount_factor=0.99)

# Hyperparameter comparison
learning_rates = [0.1, 0.9]  
gammas = [0.95, 0.99]

for lr in learning_rates:
    for gamma in gammas:
        qlearning_train_gridworld(env=env, learning_rate=lr, discount_factor=gamma)
\end{verbatim}

\subsubsection{Results Analysis}

The Q-Learning algorithm successfully learned optimal policies for the GridWorld environment:

\begin{itemize}
    \item \textbf{Convergence}: The algorithm converged to optimal policies within 5000 episodes
    \item \textbf{Path Efficiency}: Trained agents consistently found paths to the goal in 6 steps
    \item \textbf{Hyperparameter Sensitivity}: Higher learning rates (0.9) and discount factors (0.99) led to faster convergence
    \item \textbf{Q-Value Analysis}: The learned Q-values correctly reflected the value of different state-action pairs
\end{itemize}

\subsubsection{Performance Metrics}

Through TensorBoard logging, we observed:

\begin{itemize}
    \item \textbf{Sample Efficiency}: Q-Learning achieved optimal performance faster than deep RL methods
    \item \textbf{Stability}: The algorithm showed consistent learning curves across multiple runs
    \item \textbf{Exploration}: Epsilon-greedy exploration effectively balanced exploration and exploitation
\end{itemize}

\subsection{Environment Validation}

\subsubsection{State Space Verification}

The environment correctly implements the 16-state space with proper state transitions and boundary handling.

\subsubsection{Action Space Validation}

All four actions (up, down, left, right) are properly implemented with appropriate boundary constraints.

\subsubsection{Reward Function Testing}

The reward function correctly assigns:
\begin{itemize}
    \item +10 for reaching the goal
    \item -1 for falling into holes
    \item 0 for all other transitions
\end{itemize}

\subsection{Conclusion}

The custom GridWorld environment successfully demonstrates:

\begin{itemize}
    \item Proper MDP formulation with clear state, action, reward, and transition definitions
    \item Correct Gymnasium API implementation enabling compatibility with RL libraries
    \item Effective Q-Learning algorithm that learns optimal policies
    \item Comprehensive hyperparameter analysis showing the impact of learning rate and discount factor
    \item Successful validation of environment properties and agent performance
\end{itemize}

The environment serves as an excellent testbed for comparing different RL algorithms and understanding the fundamentals of reinforcement learning in discrete state spaces.

}}


%%%%%%%%%%%%%%%%%%%%%%%%%%%%%%%%%%%%%%%%%%%%%%%%%


\newpage

{\fontfamily{lmss}\selectfont {\color{DarkBlue}

\section{Task 3: Pygame for RL environment [20-points]}


}}


%%%%%%%%%%%%%%%%%%%%%%%%%%%%%%%%%%%%%%%%%%%%%%%%%

\newpage

{\fontfamily{lmss}\selectfont {\color{DarkBlue}

\begin{thebibliography}{9}

\bibitem{SuttonBarto}
R. Sutton and A. Barto, \textit{Reinforcement Learning: An Introduction}, 2nd Edition, 2020. Available online: \url{http://incompleteideas.net/book/the-book-2nd.html}


\bibitem{StableBaselines3}
A. Raffin et al., "Stable Baselines3: Reliable Reinforcement Learning Implementations," GitHub Repository, 2020. Available: \url{https://github.com/DLR-RM/stable-baselines3}.

\bibitem{Gymnasium}
Gymnasium Documentation. Available: \url{https://gymnasium.farama.org/}.

\bibitem{Pygame}
Pygame Documentation. Available: \url{https://www.pygame.org/docs/}.

\bibitem{CS285}
CS 285: Deep Reinforcement Learning, UC Berkeley, Pieter Abbeel. Course material available: \url{http://rail.eecs.berkeley.edu/deeprlcourse/}.

\bibitem{Freepik}
\href{https://www.freepik.com/free-vector/cute-artificial-intelligence-robot-isometric-icon_16717130.htm}{Cover image designed by freepik}

\end{thebibliography}


}}

%%%%%%%%%%%%%%%%%%%%%%%%%%%%%%%%%%%%%%%%%%%%%%%%%

\end{document}